% ==============================================================================
%  4.7  Validación de la capacidad predictiva
% ==============================================================================
\section{Validación de la capacidad predictiva}
\label{sec:t4-capacidad-predictiva}

La \textbf{capacidad predictiva} de un modelo es su capacidad para realizar predicciones válidas
en individuos distintos de aquellos con los que se ha construido.

\paragraph{Validación externa.} La mejor validación de un modelo es la que utiliza una
\textbf{muestra externa}, es decir, la que comprueba cómo funciona el modelo con datos nuevos que
no se han utilizado para ajustarlo. Con frecuencia esto no es posible.

\paragraph{Validación interna.} Si no se dispone de una muestra externa, se utilizan las
observaciones disponibles. Cuando el conjunto de datos es grande, una solución intermedia es
repartirlo aleatoriamente en dos grupos: una \textbf{muestra de ajuste} o de entrenamiento
(\emph{training data set}), con la que se ajusta el modelo, y una \textbf{muestra de validación}
(\emph{test data set}), con la que se evalúa. Cuando no es posible dividir la muestra de esta
forma, se recurre a métodos iterativos, que se describen a continuación.

\subsection{Validación cruzada dejando uno fuera (LOOCV)}

La validación cruzada dejando uno fuera (\emph{Leave One Out Cross-Validation}, LOOCV) es un
método iterativo:
\begin{itemize}
  \item el conjunto de entrenamiento lo forman todas las observaciones disponibles excepto una;
  \item la observación excluida es la muestra de validación;
  \item el proceso se repite tantas veces como observaciones haya, excluyendo en cada iteración
    una observación distinta, ajustando el modelo con el resto y calculando el error de
    predicción en la observación excluida.
\end{itemize}
En la iteración $i$-ésima, el error de predicción es el residuo eliminado
$e_{(i)}=y_i-\est{Y}_{i(i)}$ (\cref{sec:t4-eliminados}).

\subsection{Validación cruzada en \texorpdfstring{$k$}{k} grupos (\emph{k-fold})}

La validación cruzada en $k$ grupos (\emph{K-Fold Cross-Validation}) también es un método
iterativo:
\begin{itemize}
  \item los datos se dividen aleatoriamente en $k$ grupos de aproximadamente el mismo tamaño;
  \item $k-1$ grupos se emplean para entrenar el modelo, y el grupo restante como muestra de
    validación;
  \item el proceso se repite $k$ veces, utilizando en cada iteración un grupo distinto como
    muestra de validación.
\end{itemize}
El método LOOCV es el caso particular $k=n$.

\subsection{Bootstrap}

\begin{itemize}
  \item Es un método iterativo.
  \item Una \textbf{muestra bootstrap} es una muestra obtenida a partir de la muestra original
    mediante muestreo aleatorio con reposición, y del mismo tamaño que la muestra original.
  \item Las observaciones de la muestra original no seleccionadas en una muestra bootstrap se
    denominan \emph{out-of-bag} (OOB).
  \item En cada iteración se genera una nueva muestra bootstrap, se ajusta el modelo con ella y se
    evalúa con las observaciones \emph{out-of-bag}. La estimación del error de predicción en esa
    iteración es la media de los errores de predicción en las observaciones \emph{out-of-bag}.
  \item El proceso se repite $B$ veces y se calcula la media de las $B$ estimaciones del error de
    predicción.
\end{itemize}

\begin{ejemplo}[Muestras bootstrap]
  A partir de una muestra original $\{x_1,\ldots,x_{10}\}$ se generan tres muestras bootstrap
  (conjuntos de entrenamiento), cada una con sus observaciones \emph{out-of-bag} (conjuntos de
  validación):
  \begin{center}
    \renewcommand{\arraystretch}{1.2}
    \begin{tabular}{lll}
      \toprule
      & Muestra bootstrap (entrenamiento) & \emph{Out-of-bag} (validación) \\
      \midrule
      1 & $x_8,x_6,x_2,x_9,x_5,x_8,x_1,x_4,x_8,x_2$ & $x_3,x_7,x_{10}$ \\
      2 & $x_{10},x_1,x_3,x_5,x_1,x_7,x_4,x_2,x_1,x_8$ & $x_6,x_9$ \\
      3 & $x_6,x_5,x_4,x_1,x_2,x_4,x_2,x_6,x_9,x_2$ & $x_3,x_7,x_8,x_{10}$ \\
      \bottomrule
    \end{tabular}
  \end{center}
\end{ejemplo}
